\chapter{Foundations: forges, the archive, and the data model}
\label{ch:foundations}

Before we can name a piece of code, we need somewhere that reliably keeps it. This chapter
builds that ground: why software forges are not archives, what a universal archive of source
code collects and how, and the content-addressed data model --- the Merkle DAG --- that makes
the identifiers of the next chapter possible.

\takeaway{A forge is where code is developed; an archive is where it is kept.}

\section{Forges are not archives}
% TODO draft from: forges-not-archives.org::#categories,#oldapproaches,#evidence.
% Link-rot evidence: Google Code/Gitorious (2015), Bitbucket (2020), URL half-life ~4y.

\section{The universal archive}
% TODO draft from: swh-as-infrastructure.org::#oneslide; under-the-hood-pictures.org::#automation.
% What Software Heritage collects and how (listers + loaders), at what scale; non-profit,
% built for the long term. Refs: SwhCACM2018.

\section{The Merkle DAG data model}
% TODO draft from: under-the-hood-pictures.org::#automation; vcs-history.org::#dvcs-to-merkle.
% Content-addressed, deduplicated, history-preserving --- this is what makes intrinsic
% identification possible, and sets up Chapter~\ref{ch:identifiers}.
% Refs: pietri2019graph, SWHecosystems2023.
Software Heritage archives source code with its full development history in a uniform Merkle
DAG~\autocite{ICMS2020,SWHecosystems2023} --- the property that the next chapter turns into an
identifier you can trust.
